\documentclass[11pt]{article}
\usepackage[margin=25mm]{geometry}
\usepackage{amsmath,amssymb,amsthm,hyperref}
\newtheorem{theorem}{Theorem}
\newtheorem{lemma}{Lemma}
\title{Two short scalar constraints at cutoff five}
\author{Research note}\date{7 October 2026}
\begin{document}\maketitle
\noindent\textbf{Status.} These are proved necessary conditions for a
maximal corner. They strengthen the finite cover at the requested cutoff
$5$. They do not prove global Fourier confinement or the final approximate
alternatives $p\approx0,5,p_*$. No smaller cutoff is substituted for $5$.

Let $E$ be a corner of an order-six complex Hadamard, and use the
reconstruction invariants
\[
 p=X+Y+Z,\quad J=2\operatorname{Re}(sr\bar t),\quad
 W=XY+XZ+YZ-XYZ.
\]
In the remainder of this note, $t$ denotes the scalar $9-p$, not a Gram
entry. Put
\[
 c=W-t(t+9),\qquad
 R_\pm=4(3t\pm J)^2(t\pm J)-\bigl(c\mp(t+6)J\bigr)^2.
\]
The same-row complementary corner has the three off-diagonal Gram entries
negated. Hence its invariant is $R_-$, while $R(E)=R_+$.
For physical corners, both invariants are nonnegative and
$p+|J|\le9$. These are the complementary permanent/discriminant
identities of the reconstruction.

Choose $E$ to attain $m=M(H)>0$. Then
\[
 0\le R_-\le R_+=m,\quad
 U=\frac{R_++R_-}{2},\quad V=\frac{R_+-R_-}{2},\quad
 0\le V\le U,\quad U+V=m.
\]
In particular, if $m<5$, then $U<5$ and $0\le V<5/2$.

\begin{lemma}[A small exact identity]\label{identity}
Writing $z=J^2$ and $z_0=t^3(4-t)/4$, we have
\begin{equation}\label{main}
 4z(z-z_0)(z+36p)
 =JV(2z+30t^2)-Uz(t+6)^2-\frac{V^2}{4}.
\end{equation}
\end{lemma}
\begin{proof}
Expanding only the two complementary discriminants gives
\[
 U=36t^3-c^2-(t^2-16t+36)J^2,\qquad
 V=2J\bigl(2J^2+30t^2+c(t+6)\bigr).
\]
Substitute these into \eqref{main} and cancel. This is a polynomial
identity, including at $J=0$; there is no division or limiting argument.
The script \texttt{complementary\_scalar\_rule.py} verifies the expansion
over the integers with a common denominator four.
\end{proof}

\begin{theorem}[Negative branch]
If the maximal corner has $J<0$, then
\[
 \boxed{\quad 5<p<9,\qquad 4J^2<(9-p)^3(p-5).\quad}
\]
This conclusion does not require $m<5$.
\end{theorem}
\begin{proof}
The right side of \eqref{main} is strictly negative: its first and last
terms are nonpositive, and $U>0$, $z>0$, $t\ge |J|>0$.
As $z+36p>0$, it follows that $0<z<z_0$.
Consequently $0<t<4$, which gives the claimed interval for $p$, and
$4z<t^3(4-t)$ gives the second assertion.
\end{proof}

\begin{theorem}[Positive branch at cutoff five]
Suppose $m<5$, $J>0$, and $p\le7$. Set $t=9-p$. Then
\begin{equation}\label{positive}
 \boxed{\quad
 5\bigl(2J^2+30t^2-J(t+6)^2\bigr)
 +2J(J^2+36p)\bigl(t^3(4-t)-4J^2\bigr)\ge0.
 \quad}
\end{equation}
\end{theorem}
\begin{proof}
Here $2\le t\le9$ and $0<J\le t$. Put
$D=2J^2+30t^2-J(t+6)^2$. The factorization
\[
 \frac D J
 =\frac{2(J-t)(J-15t)}J+(t-2)(18-t)\ge0
\]
shows that $D\ge0$. Divide \eqref{main} by $J>0$, use $U\ge V$,
and discard the nonpositive last term. The result is
\[
 4J(z-z_0)(z+36p)\le VD\le\frac52D.
\]
Rearranging proves \eqref{positive}.
\end{proof}

\paragraph{What this says about stability.}
At $R_+=R_-=0$, \eqref{main} gives the critical alternatives
$J=0$ or $J^2=t^3(4-t)/4$, since $z+36p>0$ unless $p=J=0$.
For positive discriminants it gives an exact, explicitly bounded error
term, with $U<5$ and $V<5/2$. The two theorems retain useful consequences
of that error without dividing by a potentially small $J$.
Obtaining the final alternatives for $p$ still requires the additional
strip constraints; the complementary-corner identity alone does not
provide them.

\paragraph{Transparent finite checks.}
For the negative branch, the verifier encloses $J$, $p-5$, and
$(9-p)^3(p-5)-4J^2$ on each seed box using the existing certified
Taylor--Bernstein kernel. For the positive branch, it first encloses $p,J$,
then bounds the displayed two-variable polynomial using all 36 tensor
Bernstein coefficients, each evaluated with outward Arb arithmetic.
Floating-point calculations only propose which boxes to check.
Every accepted exclusion is a strict outward inequality, and all accepted
bounds are checked to be finite. The scripts
\texttt{refine\_complementary\_scalar.py} and
\texttt{refine\_positive\_scalar.py} preserve the inherited partition
exactly and retain every unresolved box.
\end{document}
